\section{Шпаргалка: Lua и Python}
\label{sec:cheatsheet}

Одно и то же действие на двух языках. Lua ждёт событий в \texttt{callback}, Python опрашивает автопилот в цикле.

\begin{center}
\footnotesize
\renewcommand{\arraystretch}{1.2}
\begin{tabularx}{\linewidth}{@{}>{\raggedright\arraybackslash}p{0.18\linewidth} >{\raggedright\arraybackslash}X >{\raggedright\arraybackslash}X@{}}
\toprule
\textbf{Действие} & \textbf{Lua} & \textbf{Python} \\
\midrule
Моторы          & \texttt{ap.push(Ev.MCE\_PREFLIGHT)} & \texttt{pioneer.arm()} \\
Взлёт           & \texttt{ap.push(Ev.MCE\_TAKEOFF)} & \texttt{pioneer.takeoff()} \\
Точка           & \texttt{ap.goToLocalPoint(x, y, z)} & \texttt{pioneer.go\_to\_local\_point(x=, y=, z=, yaw=)} \\
Ждать точку     & \texttt{if event == Ev.POINT\_REACHED} & \texttt{while not pioneer.point\_reached(): time.sleep(0.1)} \\
Курс            & \texttt{ap.updateYaw(рад)} & \texttt{yaw=} в \texttt{go\_to\_local\_point} \\
Посадка         & \texttt{ap.push(Ev.MCE\_LANDING)} & \texttt{pioneer.land()} \\
Выключить моторы & \texttt{ap.push(Ev.ENGINES\_DISARM)} & \texttt{pioneer.disarm()} \\
Пауза           & \texttt{Timer.callLater(t, f)} & \texttt{time.sleep(t)} \\
Периодически    & \texttt{Timer.new(p, f):start()} & цикл \texttt{while} с \texttt{time.sleep(p)} \\
Светодиоды      & \texttt{leds:set(i, r, g, b)}, 0…1 & \texttt{pioneer.led\_control(led\_id, r, g, b)}, 0…255 \\
Позиция         & \texttt{Sensors.lpsPosition()} & \texttt{pioneer.get\_local\_position\_lps()} \\
Дальномер       & \texttt{Sensors.range()} & \texttt{pioneer.get\_dist\_sensor\_data()} \\
Батарея         & \texttt{Sensors.battery()} & \texttt{pioneer.get\_battery\_status()} \\
Кадр камеры     & \texttt{camera.requestMakeShot()} & \texttt{Camera().get\_cv\_frame()} \\
Вывод           & \texttt{print(...)} & \texttt{print(...)} \\
Завершение      & --- & \texttt{pioneer.close\_connection()} \\
\bottomrule
\end{tabularx}
\end{center}

\begin{infobox}[Три главные мысли]
\begin{itemize}
    \item Дрон --- автопилот: вы отдаёте команду и ждёте события, что она выполнена.
    \item В Lua ждут через \texttt{callback(event)}, в Python --- опросом \texttt{while not ...}.
    \item Координаты --- метры от старта, Z вверх, курс в радианах; их показывает \key{Ctrl}+щелчок в симуляторе.
\end{itemize}
\end{infobox}
